\section*{الفصل \ref{ch:b1:geological-maps} --- الخرائط الجيولوجية}

\begin{solution}{exo:b1:geological-maps:1}
\omterm{def:b1:geological-maps:attitude}{الامتداد}: سمت الخطوط الأفقية في المستوي؛ \omterm{def:b1:geological-maps:attitude}{الانحدار}: الزاوية بين المستوي والأفق، مقيسةً عموديًا على \omterm{def:b1:geological-maps:attitude}{الامتداد}؛ \omterm{def:b1:geological-maps:attitude}{اتجاه الانحدار}: سمت الخط الهابط الأشدّ انحدارًا. وبالنسبة إلى $135/40$، \omterm{def:b1:geological-maps:attitude}{فالامتداد} وفق قاعدة اليد اليمنى هو $135 - 90 = \ang{045}$. الرمز: خط امتداد موجّه
\mbox{\textdir TLT N\,45\,E}، وعليه شرطة على جانبه الجنوبي الشرقي (نحو \ang{135}) مكتوب عليها 40.
\end{solution}

\begin{solution}{exo:b1:geological-maps:2}
$z = 500 - \tan\ang{30}\,(x\sin\ang{90} + y\cos\ang{90}) = 500 - 0.577\,x$.
عند $(1000, 0)$: $500 - 577 = \qty{-77}{m}$. وعند $(0, 1000)$: \qty{500}{m}، لأن المستوي أفقي على طول امتداده الشمالي الجنوبي.
\end{solution}

\begin{solution}{exo:b1:geological-maps:3}
الأثر الذي يتبع كنتورًا هو \omterm{def:b1:geological-maps:geological-map}{تماس} أفقي، عند \qty{350}{m}. والأثر الذي يمتد مستقيمًا عبر التلال والأودية ينتمي إلى \omterm{def:b1:geological-maps:geological-map}{تماس} شاقولي.
\end{solution}

\begin{solution}{exo:b1:geological-maps:4}
على طول النهر، تُظهر لبّ \omterm{def:b1:geological-maps:fold}{الطية المحدبة} الشيست والغرانيت (تحت الحجر الرملي)؛ ولبّ \omterm{def:b1:geological-maps:fold}{الطية المقعرة} يشغله الطباشير. فأقدم الوحدات تقع أدنى التتابع: حيث تنحني الطبقات نحو الأعلى ويقطعها التعرية أعمق ما يكون، تظهر أقدم الوحدات؛ وحيث تنخفض تُحفظ أحدث الوحدات في الوسط.
\end{solution}

\begin{solution}{exo:b1:geological-maps:5}
أول نقطتين على الارتفاع نفسه: فالخط $y = 0$ خط امتداد، إذن \omterm{def:b1:geological-maps:attitude}{الامتداد} شرقي غربي. ويكون \omterm{def:b1:geological-maps:geological-map}{التماس} أدنى بمقدار \qty{100}{m} عند \qty{500}{m} نحو الشمال: فهو ينحدر نحو الشمال، $\tan\delta = 100/500 = 0.2$، و$\delta \approx \ang{11.3}$؛ \omterm{def:b1:geological-maps:attitude}{واتجاه الانحدار} \ang{000}، \omterm{def:b1:geological-maps:attitude}{والامتداد} \ang{270} وفق قاعدة اليد اليمنى. وبالنظرية: $\overrightarrow{P_1P_2} = (1000, 0, 0)$،
$\overrightarrow{P_1P_3} = (0, 500, -100)$، و$\vec n = (0, 100\,000,
500\,000)$، و$\cos\delta = 500\,000/509\,902 \approx 0.981$، و$\delta \approx
\ang{11.3}$، و$(n_x, n_y) = (0, 100\,000)$ يشير نحو الشمال: النتيجة نفسها.
\end{solution}

\begin{solution}{exo:b1:geological-maps:6}
الاتجاه نحو المصب هو الشمال، و$s = 0.03$ (\ang{1.72}). (أ) $a = 0 < s$: الحرف V نحو المنبع (الجنوب)، على طول الكنتورات. (ب) $a = -\tan\ang{20} < 0 < s$: الحرف V نحو المنبع. (ج) $a = \tan\ang{1} \approx 0.017 < s$: الحرف V نحو المنبع. (د) $a = \tan\ang{20}
\approx 0.364 > s$: الحرف V نحو المصب (الشمال).
\end{solution}

\begin{solution}{exo:b1:geological-maps:7}
$\xi_{\text{القاع}} - \xi_{\text{الجانب}} = h/(a - s) = 80/(0.364 - 0.03)
\approx \qty{240}{m}$: يقع الأثر على الجانب نحو \qty{240}{m} أعلى المنبع من رأس الحرف.
\end{solution}

\begin{solution}{exo:b1:geological-maps:8}
يشير اتجاه إغلاق \omterm{def:b1:geological-maps:fold}{الطية المحدبة} نحو اتجاه الغوص: فمحورها يغوص نحو الشمال. ومع محور يغوص نحو الشمال، تُغلق \omterm{def:b1:geological-maps:geological-map}{انكشافات} \omterm{def:b1:geological-maps:fold}{الطية المقعرة} نحو الجنوب.
\end{solution}

\begin{solution}{exo:b1:geological-maps:9}
الجانب الشرقي الجنوبي الشرقي يُظهر وحدات أحدث (الطباشير والصخر الطيني في مواجهة وحدات أقدم): فتلك الكتلة هي التي هبطت. وتقع شرطات رمز الصدع على ذلك الجانب، جانب الهبوط.
\end{solution}

\begin{solution}{exo:b1:geological-maps:10}
يقع خط مشترك للمستويين عموديًا على الناظمين معًا، فهو إذن موازٍ للمتجهة $\vec n_1 \times \vec n_2$. بالنسبة إلى $090/30$: $\vec n_1 = (0.577, 0,
1)$؛ وبالنسبة إلى $270/30$: $\vec n_2 = (-0.577, 0, 1)$. فيكون $\vec n_1 \times \vec n_2
= (0\cdot 1 - 1\cdot 0,\ 1\cdot(-0.577) - 0.577\cdot 1,\ 0) = (0, -1.155,
0)$: خط أفقي شمالي جنوبي. فالمستوي الأول ينحدر نحو الشرق، والثاني نحو الغرب: إنهما ينحدران بعيدًا عن الخط، وهو قمة \omterm{def:b1:geological-maps:fold}{طية} محدبة محورها أفقي.
\end{solution}

\begin{solution}{exo:b1:geological-maps:11}
أثر السدّ الشاقولي هو الخط المستقيم المار بنقطة الأصل والممتد باتجاه \ang{045}، أيًّا كان الارتفاع: فعلى التلّة يقع أثره في مكان ما على الخط نفسه. أما السدّ المائل فهو المستوي $z = -\tan\ang{60}\,d$، حيث $d$ المسافة الأفقية عن خط \omterm{def:b1:geological-maps:attitude}{الامتداد} المارّ بنقطة الأصل، مأخوذة موجبة نحو الجنوب الشرقي. وعند $z = \qty{300}{m}$: $d = -300/\tan\ang{60}
\approx \qty{-173}{m}$: يقع الأثر على بُعد \qty{173}{m} شمال غرب خط \omterm{def:b1:geological-maps:attitude}{الامتداد} المارّ بنقطة الأصل، على الجانب المعاكس \omterm{def:b1:geological-maps:attitude}{لاتجاه الانحدار}.
\end{solution}

\begin{solution}{exo:b1:geological-maps:12}
\[
  \frac{\partial z}{\partial x} \approx \frac{405 - 420}{100} = -0.15, \qquad
  \frac{\partial z}{\partial y} \approx \frac{412 - 420}{100} = -0.08,
\]
إذن الميل هو $\arctan\sqrt{0.15^2 + 0.08^2} = \arctan 0.17 \approx
\ang{9.6}$. فالأرض تنحدر على طول $-\nabla z = (0.15, 0.08)$، وسمته هو
$\arctan(0.15/0.08) \approx \ang{62}$: فالخلية تواجه شرقًا-شمالًا شرقيًا.
\end{solution}

\begin{solution}{pb:b1:geological-maps:1}
\textbf{1.} الشيست، الغرانيت، الحجر الرملي، الحجر الجيري، الصخر الطيني، الطباشير، الطمي.

\textbf{2.} من الغرب إلى الشرق: الطباشير والصخر الطيني، ثم الحجر الجيري، ثم الحجر الرملي، ثم الشيست والغرانيت قرب $x = \qty{3.5}{km}$، ثم الحجر الرملي والحجر الجيري والصخر الطيني من جديد، ثم الطباشير؛ والطمي على طول القاع كله.

\textbf{3.} الحجر الرملي، والشيست والغرانيت تحته على طول النهر: أقدم الوحدات في اللبّ، وهي \omterm{def:b1:geological-maps:fold}{طية} محدبة.

\textbf{4.} الطباشير، أحدث الوحدات: \omterm{def:b1:geological-maps:fold}{طية} مقعرة.

\textbf{5.} قاعدة الحجر الرملي أعلى ما تكون على محور \omterm{def:b1:geological-maps:fold}{الطية المحدبة}، ولا يقطع التعرية أسفلها إلا على طول النهر.

\textbf{6.} عدم توافق: يستقر الحجر الرملي على وحدتين أقدم مختلفتين، الشيست والغرانيت، ويقطع خط محيط الغرانيت.

\textbf{7.} $435 - 261 = 174$, $432 - 428 = 4$, $512 - 586 = \qty{-74}{m}$.

\textbf{8.} $\overrightarrow{E_1E_2} = (600, 1200, -170)$,
$\overrightarrow{E_1E_3} = (1100, 400, -248)$.

\textbf{9.} الجداء المتجهي هو
\[
  (-229\,600,\ -38\,200,\ -1\,080\,000);
\]
موجّهًا نحو الأعلى، $\vec n_E = (229\,600,\ 38\,200,\ 1\,080\,000)$.

\textbf{10.} $\|\vec n_E\| \approx 1\,104\,800$; $\cos\delta \approx 0.9776$,
$\delta \approx \ang{12.2}$.

\textbf{11.} \omterm{def:b1:spherical-astronomy:horizontal}{السمت} للمتجهة $(229\,600, 38\,200)$: $\arctan(229\,600/38\,200)
\approx \ang{80.6}$؛ \omterm{def:b1:geological-maps:attitude}{والوضعية} $081/12$، \omterm{def:b1:geological-maps:attitude}{والامتداد} $\approx \ang{351}$.

\textbf{12.} نعم: للرمز على الجناح الشرقي خط امتداد قرب الشمالي الجنوبي، وشرطة انحدار نحو الشرق، مكتوب عليها 12.

\textbf{13.} $524 - 669 = -145$, $421 - 459 = -38$, $319 - 209 =
\qty{110}{m}$.

\textbf{14.} $\overrightarrow{W_1W_2} = (700, 1200, 107)$ و
$\overrightarrow{W_1W_3} = (1100, -600, 255)$. وجداؤهما المتجهي، موجّهًا نحو الأعلى، هو
\[
  (370\,200,\ -60\,800,\ -1\,740\,000) \quad\longrightarrow\quad
  \vec n_W = (-370\,200,\ 60\,800,\ 1\,740\,000).
\]

\textbf{15.} $\cos\delta = 1\,740\,000/\|\vec n_W\| \approx 0.9776$، $\delta
\approx \ang{12.2}$؛ \omterm{def:b1:geological-maps:attitude}{واتجاه الانحدار}: سمت $(-370\,200, 60\,800)$،
$\approx \ang{279.3}$. \omterm{def:b1:geological-maps:attitude}{والوضعية} $279/12$، \omterm{def:b1:geological-maps:attitude}{والامتداد} $\approx \ang{189}$: ينحدر الجناح الغربي نحو الغرب-الشمال الغربي، والجناح الشرقي نحو الشرق-الشمال الشرقي.

\textbf{16.} $\vec n_E \times \vec n_W = (804 \times 10^6,\ -799\,320 \times
10^6,\ 28\,101 \times 10^6)$ (مقرّبًا).

\textbf{17.} اتجاه النزول هو $(-804, 799\,320, -28\,101) \times
10^6$: الاتجاه $\arctan(-804/799\,320) \approx \ang{0}$، أي نحو الشمال؛ والغوص
$\arctan\bigl(28\,101/\sqrt{804^2 + 799\,320^2}\bigr) \approx \ang{2.0}$. فالمحور يغوص \ang{2} نحو الشمال.

\textbf{18.} $8000 \times \tan\ang{2.0} \approx \qty{280}{m}$ أدنى، أي نحو
\qty{-50}{m}. ولبّ \omterm{def:b1:geological-maps:fold}{الطية المحدبة} أعلى ما يكون قرب النهر، وينخفض نحو الشمال، فينغلق \omterm{def:b1:geological-maps:geological-map}{انكشافها} نحو الشمال كما تُظهر الخريطة.

\textbf{19.} $\overrightarrow{F_1F_2} = (1627, 2500, -276)$ و
$\overrightarrow{F_1F_3} = (3124, 5500, 77)$ تعطيان الناظم نحو الأعلى
\[
  (1\,710\,500,\ -987\,503,\ 1\,138\,500), \qquad
  \delta = \arccos\frac{1\,138\,500}{2\,279\,700} \approx \ang{60},
\]
وسمت \omterm{def:b1:geological-maps:attitude}{اتجاه الانحدار} $\approx \ang{120}$. \omterm{def:b1:geological-maps:attitude}{فالوضعية} $120/60$، \omterm{def:b1:geological-maps:attitude}{والامتداد}
\ang{030}.

\textbf{20.} ينحدر الصدع نحو الشرق-الجنوب الشرقي، فهذا الجانب هو الجدار العلوي، أي الكتلة الواقعة فوق مستوي الصدع المائل؛ وقد هبطت (وحدات أحدث على السطح): صدع عادي.

\textbf{21.} ينحدر مستوي الصدع نحو الشرق-الجنوب الشرقي، أي نحو المصب، وبانحدار أشدّ بكثير من انحدار قاع النهر: فتلتقي الأرض الأخفض بالمستوي عند نقطة أبعد على طول \omterm{def:b1:geological-maps:attitude}{الانحدار}، فينحرف الأثر نحو الشرق-الجنوب الشرقي في الوادي --- قاعدة V مطبّقة على مستوي الصدع.

\textbf{22.} اقتحام الغرانيت؛ التعرية قبل الحجر الرملي؛ ترسيب الطباشير؛ الطيّ؛ الصدع؛ ترسيب الطمي.

\textbf{23.} بعد 90 مليون سنة مضت (الطباشير مطوية) وقبل 2.58 مليون سنة مضت (الطمي غير مصدوع)؛ فالطيّ جاء أولًا، والصدع بعده.

\textbf{24.} ينزل الصدع \qty{1405}{m} من F1 (\qty{405}{m}) إلى
\qty{-1000}{m}؛ وعند انحدار \ang{60} يتطلب ذلك بعدًا أفقيًا مقداره
$1405/\tan\ang{60} \approx \qty{811}{m}$ على طول \omterm{def:b1:geological-maps:attitude}{الانحدار}. ويجب حفر البئر على بُعد نحو \textbf{\qty[text-series-to-math=true]{810}{m} شرق-جنوب شرق F1} (\omterm{def:b1:spherical-astronomy:horizontal}{السمت} \ang{120}).
\end{solution}
